% !TeX program = xelatex
% ============================================================
% pmi/pmi.tex — Программа и методика испытаний
% ГОСТ 19.301-79 «Программа и методика испытаний.
% Требования к содержанию и оформлению»
% ============================================================
\documentclass[12pt,a4paper]{extarticle}

\input{../common/preamble}

\newcommand{\docname}{Программа и методика испытаний}
\newcommand{\doccode}{\hseDocCodeBase\ 51 01-1}
\newcommand{\doccodelu}{\hseDocCodeBase\ 51 01-1-ЛУ}

\input{../common/commands}
\input{../common/styles}

\begin{document}
\sloppy

% ── Титульный лист ──────────────────────────────────────────
\input{../common/title_template}


% ============================================================
% АННОТАЦИЯ
% ============================================================
\section*{АННОТАЦИЯ}
\addcontentsline{toc}{section}{АННОТАЦИЯ}

Настоящая Программа и методика испытаний на разработку программы «\hseProjectName» содержит следующие разделы: «Объект испытаний», «Цель испытаний», «Требования к~программе», «Требования к~программной документации», «Средства и~порядок испытаний», «Методы испытаний».

В разделе «Объект испытаний» указаны наименование, область применения и~обозначение испытуемой программы.

В разделе «Цель испытаний» указана цель проведения испытаний.

Раздел «Требования к~программе» содержит перечень требований из~технического задания, подлежащих проверке во~время испытаний.

Раздел «Требования к~программной документации» содержит состав программной документации, предъявляемой на~испытания.

Раздел «Средства и~порядок испытаний» содержит информацию о~технических и~программных средствах, используемых во~время испытаний, а~также порядок их проведения.

В разделе «Методы испытаний» приведено описание используемых методов и~конкретных проверок.

% Подсказка: перечислите, что приведено в приложениях вашей ПМИ.
В приложениях приведены перечень терминов, матрица трассировки требований к~автоматизированным тестам и~перечень подтверждающих материалов испытаний.

\clearpage

% ============================================================
% СПИСОК ИСПОЛЬЗУЕМЫХ ГОСТ
% ============================================================
\noindent
Настоящий документ разработан в~соответствии с~требованиями:

\begin{enumerate}[label=\arabic*), leftmargin=1.25cm]
\item ГОСТ 19.101-77 Виды программ и~программных документов\cite{gost19101}.
\item ГОСТ 19.103-77 Обозначения программ и~программных документов\cite{gost19103}.
\item ГОСТ 19.104-78 Основные надписи\cite{gost19104}.
\item ГОСТ 19.105-78 Общие требования к~программным документам\cite{gost19105}.
\item ГОСТ 19.106-78 Требования к~программным документам, выполненным
      печатным способом\cite{gost19106}.
\item ГОСТ 19.301-79 Программа и~методика испытаний. Требования
      к~содержанию и~оформлению\cite{gost19301}.
\end{enumerate}

Изменения к~данному документу оформляются согласно ГОСТ\,19.604-78\cite{gost19604}.

\clearpage

% ============================================================
% СОДЕРЖАНИЕ
% ============================================================
\hypersetup{linkcolor=black}
\tableofcontents
\hypersetup{linkcolor=blue}
\thispagestyle{fancy}
\clearpage

% ============================================================
% 1. ОБЪЕКТ ИСПЫТАНИЙ
% ============================================================
\section{ОБЪЕКТ ИСПЫТАНИЙ}

\subsection{Наименование испытуемой программы}

Наименование программы~--- «\projectname».

% Английское наименование подставляется из настроек проекта (поле
% «Название работы на английском»).
Наименование программы на английском языке~--- «\hseProjectNameEn».

\subsection{Область применения испытуемой программы}

% Подсказка: кратко опишите назначение и область применения программы,
% а также какой контур системы рассматривается в рамках испытаний.

Программа «\projectname» представляет собой \hseFill{краткое описание назначения системы}. Система охватывает \hseFill{основные домены и подсистемы}.

В рамках настоящего документа рассматривается \hseFill{испытываемый контур системы}.

\subsection{Обозначение испытуемой программы}

Наименование темы разработки~--- «\projectname».

% Краткое наименование подставляется из настроек проекта (поле
% «Краткое наименование программы»).
Краткое наименование программы~--- «\hseProjectNameShort».

\clearpage

% ============================================================
% 2. ЦЕЛЬ ИСПЫТАНИЙ
% ============================================================
\section{ЦЕЛЬ ИСПЫТАНИЙ}

Цель проведения испытаний~--- проверка соответствия характеристик разработанной программы функциональным и~иным требованиям, изложенным в~программной документации «Техническое задание».

\clearpage

% ============================================================
% 3. ТРЕБОВАНИЯ К ПРОГРАММЕ
% ============================================================
\section{ТРЕБОВАНИЯ К ПРОГРАММЕ}

\subsection{Требования к функциональным характеристикам}

\subsubsection{Состав выполняемых функций}

% Подсказка: идентификаторы (ТЗ-Ф-01 и т.д.) связывают требования
% с тестами в панели «Трассировка требований».

Проверке подлежат функциональные характеристики программы, перечисленные в~техническом задании в~таблице требований. Каждому функциональному требованию с~идентификатором вида \texttt{ТЗ-Ф-XX} сопоставляется автоматизированный тест, тестовый набор или иная проверка.

\input{../technical_specification/requirements_table}

\clearpage

\subsubsection{Организация входных данных}

% Подсказка: перечислите проверяемые входные данные и форматы
% вашей системы (структура запросов, валидация, ограничения).
\begin{hseExample}
При испытаниях проверяется выполнение требований технического задания к~организации входных данных. Проверке подлежат:

\begin{itemize}[leftmargin=1.25cm]
    \item прием \hseFill{вид запросов или входных данных} с~передачей данных в~формате \hseFill{формат входных данных};
    \item валидация входящих данных на~соответствие схемам и~типам данных, указанным в~спецификации интерфейса;
    \item соблюдение ограничений на~размер \hseFill{вид ограничений на входные данные};
    \item \hseFill{другие проверки организации входных данных}.
\end{itemize}

Проверка выполняется как на~корректных входных данных, так и~на~ошибочных наборах данных, предназначенных для подтверждения работы валидации.
\end{hseExample}

\subsubsection{Организация выходных данных}

% Подсказка: перечислите проверяемые выходные данные (коды ответов,
% форматы, структура ошибок).
\begin{hseExample}
При испытаниях проверяется выполнение требований технического задания к~организации выходных данных. Проверке подлежат:

\begin{itemize}[leftmargin=1.25cm]
    \item возврат результатов в~формате \hseFill{формат выходных данных};
    \item использование \hseFill{коды состояния или статусы ответов} для успешных и~ошибочных ответов;
    \item возврат структурированного описания ошибки, достаточного для диагностики без раскрытия конфиденциальных данных;
    \item \hseFill{другие проверки организации выходных данных}.
\end{itemize}
\end{hseExample}

\subsubsection{Требования к временным характеристикам}

% Подсказка: перечислите временные показатели из ТЗ, подлежащие
% проверке (время отклика, пропускная способность).
\begin{hseExample}
При испытаниях проверяются временные характеристики, установленные в~техническом задании:

\begin{itemize}[leftmargin=1.25cm]
    \item время обработки \hseFill{вид запросов или операций} не~должно превышать \hseFill{целевое время отклика};
    \item указанные показатели должны сохраняться при штатной нагрузке до~\hseFill{целевое число пользователей};
    \item целевая интенсивность запросов должна соответствовать диапазону \hseFill{целевая интенсивность запросов}.
\end{itemize}
\end{hseExample}

\subsection{Требования к интерфейсу}

% Подсказка: опишите проверяемые требования к программному интерфейсу
% (методы API, контракты, спецификации).
\begin{hseExample}
При испытаниях проверяется выполнение требований технического задания к~программному интерфейсу. Проверке подлежат:

\begin{itemize}[leftmargin=1.25cm]
    \item реализация \hseFill{вид интерфейса программы} в~соответствии с~\hseFill{принципы и спецификации интерфейса};
    \item доступность актуальной спецификации интерфейса \hseFill{расположение спецификации};
    \item проверка доступа к~защищенным функциям на~основе \hseFill{механизм аутентификации и авторизации};
    \item единообразие структуры ответов, форматов данных и~способов именования полей.
\end{itemize}
\end{hseExample}

\subsection{Требования к маркировке и версионированию}

% Подсказка: скорректируйте список под требования вашего ТЗ,
% если они заданы.
\begin{hseExample}
При испытаниях проверяется выполнение требований технического задания к~маркировке и~версионированию компонентов системы. Проверке подлежат:

\begin{itemize}[leftmargin=1.25cm]
    \item соответствие версий программного кода спецификации \hseFill{используемая схема версионирования};
    \item наличие версионных тегов \hseFill{вид артефактов сборки};
    \item наличие метаданных сборки с~датой сборки и~идентификатором версии исходного кода;
    \item хранение исходного кода, конфигураций и~артефактов сборки в~хранилищах, поддерживающих контроль версий.
\end{itemize}
\end{hseExample}

\clearpage

% ============================================================
% 4. ТРЕБОВАНИЯ К ПРОГРАММНОЙ ДОКУМЕНТАЦИИ
% ============================================================
\section{ТРЕБОВАНИЯ К ПРОГРАММНОЙ ДОКУМЕНТАЦИИ}

\subsection{Состав программной документации, предъявляемой на испытания}

% Подсказка: перечислите документы, предъявляемые на испытания,
% с указанием соответствующих ГОСТ.

На испытания должна быть представлена документация к~программе в~следующем составе:
\begin{enumerate}[label=\arabic*), leftmargin=1.25cm]
    \item «\projectname». Техническое задание (ГОСТ 19.201-78).
    \item «\projectname». Текст программы (ГОСТ 19.401-78).
    \item «\projectname». Руководство программиста (ГОСТ 19.504-79).
    \item «\projectname». Программа и~методика испытаний
          (ГОСТ 19.301-79).
\end{enumerate}

\subsection{Специальные требования}

Документы к~программе должны быть выполнены в~соответствии с~ГОСТ 19.106-78\cite{gost19106} и~ГОСТами к~каждому виду документа (см.~п.~4.1).

% Подсказка: при необходимости укажите иные специальные требования
% к оформлению документации.
Подтверждающие материалы испытаний допускается оформлять в~виде приложений к~настоящему документу.

\clearpage

% ============================================================
% 5. СРЕДСТВА И ПОРЯДОК ИСПЫТАНИЙ
% ============================================================
\section{СРЕДСТВА И ПОРЯДОК ИСПЫТАНИЙ}

\subsection{Технические средства, используемые во время испытаний}

% Подсказка: опишите аппаратное окружение испытаний (ОС, процессор,
% память) и контуры запуска.
\begin{hseExample}
Испытания проводятся на~технических средствах со~следующими характеристиками: \hseFill{ОС, процессор, память}. Используются следующие контуры испытаний:

% Таблица набрана через \captionof (не \begin{table}[H]): плавающие
% окружения внутри жёлтого блока-образца недопустимы.
\begin{center}
\small
\captionof{table}{Контуры проведения испытаний}
\label{tab:pmi_test_environments}
\begin{tabularx}{\textwidth}{|>{\raggedright\arraybackslash}p{0.21\textwidth}|>{\raggedright\arraybackslash}p{0.29\textwidth}|X|}
\hline
\textbf{Контур} & \textbf{Состав} & \textbf{Назначение} \\
\hline
Локальный контур & \hseFill{состав локального контура} & Быстрый запуск тестов без полного разворачивания окружения. \\
\hline
\hseFill{интеграционный контур} & \hseFill{состав интеграционного контура} & Полноценные интеграционные проверки в~окружении, максимально близком к~среде эксплуатации. \\
\hline
\hseFill{приемочный стенд} & \hseFill{состав приемочного стенда} & Приемочные проверки эксплуатационных характеристик развернутой системы. \\
\hline
\end{tabularx}
\end{center}
\end{hseExample}

\subsection{Программные средства, используемые во время испытаний}

% Подсказка: перечислите программные средства и инструменты
% (среда выполнения, фреймворки тестирования, средства контроля
% качества).
\begin{hseExample}
При испытаниях используются следующие программные средства:

\begin{itemize}[leftmargin=1.25cm]
    \item среда выполнения: \hseFill{языки и среды выполнения};
    \item фреймворки тестирования: \hseFill{фреймворки тестирования};
    \item средства статического анализа: \hseFill{линтеры и анализаторы};
    \item средства контроля качества и~покрытия: \hseFill{инструменты контроля качества};
    \item инфраструктурные средства: \hseFill{СУБД, брокеры, контейнеризация}.
\end{itemize}

Основным способом приемки функциональных требований в~настоящей работе является прохождение автоматизированных тестов: входные данные, действия и~ожидаемые результаты зафиксированы в~исполняемых тестовых сценариях, а~результаты их выполнения сохраняются в~\hseFill{CI/CD или отчеты испытаний}.
\end{hseExample}

\subsection{Приемочный критерий функциональных требований}

\begin{hseExample}
Функциональное требование с~идентификатором вида \texttt{ТЗ-Ф-XX} считается принятым, если одновременно выполняются следующие условия:

\begin{enumerate}[label=\arabic*), leftmargin=1.25cm]
    \item требование присутствует в~таблице требований технического задания;
    \item для требования указан автоматизированный тест, тестовый набор или сквозной сценарий;
    \item соответствующая проверка на~проверяемой версии программы завершилась успешно;
    \item для проверки доступны подтверждающие артефакты: журнал выполнения, отчет о~тестировании или отчет покрытия, если они предусмотрены данным контуром.
\end{enumerate}

Приемочная комиссия подтверждает выполнение функциональных требований одним из двух способов:

% Подсказка: уточните способы подтверждения под ваш проект.
\begin{enumerate}[label=\arabic*), leftmargin=1.25cm]
    \item запросить доступ к~\hseFill{репозиторий проекта}, открыть результаты проверок проверяемой версии и~убедиться в~их успешном завершении;
    \item выполнить тесты локально на~согласованной версии исходного кода командой \hseFill{команда запуска тестов}.
\end{enumerate}

Локальный запуск используется как альтернативный способ подтверждения при отсутствии доступа к~CI/CD или при необходимости независимого воспроизведения результата. При наличии расхождений между локальным запуском и~CI/CD приемочным результатом считается запуск на~согласованной версии в~изолированном CI/CD-контуре, так как он фиксирует версию кода, конфигурацию, окружение и~артефакты выполнения.
\end{hseExample}

\subsection{Порядок проведения испытаний}

% Подсказка: скорректируйте последовательность этапов под ваш проект.
\begin{hseExample}
Испытания проводятся в~следующем порядке:

\begin{enumerate}[label=\arabic*), leftmargin=1.25cm]
    \item \textbf{Подготовительный этап}: визуальная проверка состава программной документации, подготовка тестовой среды согласно разделам 5.1--5.2 и~развертывание программы согласно руководству программиста.
    \item \textbf{Функциональная приемка}: проверка требований технического задания по~результатам проверок согласно разделу~\ref{sec:methods}.
    \item \textbf{Статический контроль качества}: запуск и~анализ результатов \hseFill{линтеры и проверки качества}.
    \item \textbf{Испытания уровня модулей и~сервисов}: выполнение \hseFill{команды или задания тестирования} и~анализ отчетов о~тестировании и~покрытии.
    \item \textbf{Сквозные испытания}: выполнение сквозных сценариев в~\hseFill{контур сквозных испытаний}.
    \item \textbf{Фиксация результатов}: сравнение полученных результатов с~ожидаемыми и~оформление заключения о~прохождении испытаний.
\end{enumerate}
\end{hseExample}

\subsection{Количественные и качественные характеристики, подлежащие оценке}

% Подсказка: скорректируйте перечень оцениваемых характеристик
% под ваш проект (наблюдаемость, нагрузка и т.д.).
\begin{hseExample}
В ходе испытаний оцениваются следующие характеристики:

\begin{itemize}[leftmargin=1.25cm]
    \item процент успешно выполненных тестов и~проверок;
    \item отсутствие непредусмотренных отказов, падений и~неконсистентных состояний данных;
    \item корректность \hseFill{ключевые проверяемые аспекты программы};
    \item наличие и~качество подтверждающих артефактов испытаний;
    \item \hseFill{дополнительные характеристики из ТЗ}.
\end{itemize}
\end{hseExample}

\subsection{Условия проведения испытаний}

% Подсказка: укажите необходимые условия (доступы, версии исходного
% кода, окружение, секреты).
\begin{hseExample}
Испытания должны проводиться при соблюдении следующих условий:

\begin{itemize}[leftmargin=1.25cm]
    \item доступ к~\hseFill{репозиторий проекта} и~артефактам проверок;
    \item известный идентификатор проверяемой версии: идентификатор коммита, тег версии или ссылка на~запрос слияния;
    \item наличие \hseFill{среда выполнения программы} и~сетевого доступа к~необходимым ресурсам;
    \item подготовленные переменные окружения, секреты и~ключи тестовых контуров;
    \item запуск испытаний на~согласованной версии исходного кода и~тестовой конфигурации.
\end{itemize}
\end{hseExample}

\clearpage

% ============================================================
% 6. МЕТОДЫ ИСПЫТАНИЙ
% ============================================================
\section{МЕТОДЫ ИСПЫТАНИЙ}
\label{sec:methods}

\begin{hseExample}
Испытания выполняются сочетанием визуального контроля документации, автоматизированных тестовых наборов, контрольных проверок качества и~приемочных проверок. Для каждого функционального требования используется трассировка до~конкретного теста, тестового набора или проверки из~таблицы методов испытаний.
\end{hseExample}

% Подсказка: для каждого требования опишите метод проверки, ожидаемый
% и фактический результаты. Заполните таблицу методов испытаний.
% Таблица методов испытаний остаётся вне жёлтого блока-образца:
% longtable внутри hseExample недопустим.

\begin{longtable}{|p{0.18\textwidth}|p{0.22\textwidth}|p{0.25\textwidth}|p{0.22\textwidth}|}
\caption{Методы испытаний}
\label{tab:test-methods}\\
\hline
\textbf{Проверяемое требование} &
\textbf{Метод проверки} &
\textbf{Ожидаемый результат} &
\textbf{Фактический результат} \\
\hline
\endfirsthead
\hline
\textbf{Проверяемое требование} &
\textbf{Метод проверки} &
\textbf{Ожидаемый результат} &
\textbf{Фактический результат} \\
\hline
\endhead
\hline
\endfoot
\hline
\endlastfoot
\reqref{ТЗ-Ф-01} &
\hseFill{метод проверки} &
\hseFill{ожидаемый результат} &
\hseFill{фактический результат} \\
\hline
\end{longtable}

\subsection{Испытание выполнения требований к программной документации}

Проверка выполняется визуальным методом. Для каждого документа из~раздела~4.1 проверяются:

\begin{itemize}[leftmargin=1.25cm]
    \item наличие документа в~составе комплекта;
    \item соответствие обозначению, титульному листу и~структуре разделов;
    \item отсутствие ссылок на~устаревшие компоненты, интерфейсы и~тестовые контуры;
    \item согласованность настоящего документа с~техническим заданием, текстом программы и~руководством программиста.
\end{itemize}

\subsection{Испытание функциональных требований}

% Подсказка: скорректируйте шаги под ваш CI/CD или способ запуска тестов.
\begin{hseExample}
Проверка выполняется автоматизированным методом по~таблице трассировки требований к~тестам. Испытатель выполняет следующие действия:

\begin{enumerate}[label=\arabic*), leftmargin=1.25cm]
    \item открыть \hseFill{репозиторий или CI/CD проекта};
    \item выбрать проверку, соответствующую проверяемому коммиту или тегу версии;
    \item убедиться, что все задания тестирования завершены успешно;
    \item для каждого требования \texttt{ТЗ-Ф-XX} сопоставить указанный тест или тестовый набор с~успешно выполненной проверкой;
    \item сохранить или приложить к~материалам испытаний снимки экранов и~отчеты о~выполнении.
\end{enumerate}

Испытание считается успешным, если все функциональные требования имеют трассировку до~автоматизированной проверки, а~соответствующие проверки завершены без неразобранных падений тестов. Если отдельный тест был повторно запущен из-за нестабильности окружения, в~материалах испытаний фиксируется итоговый успешный запуск и~ссылка на~журнал повторного выполнения.
\end{hseExample}

\subsection{Испытание статического контроля качества и проверок безопасности}

% Подсказка: перечислите используемые линтеры, анализаторы качества
% и проверки безопасности; удалите подраздел, если они не применяются.
\begin{hseExample}
Проверка выполняется автоматизированным методом. Используются следующие группы проверок:

\begin{itemize}[leftmargin=1.25cm]
    \item \hseFill{линтеры и форматтеры} для анализа форматирования, линтинга и~базовых правил качества;
    \item \hseFill{платформа анализа качества} для построения аналитики качества и~агрегирования покрытия тестами;
    \item \hseFill{средства проверки безопасности} для поиска уязвимостей зависимостей и~секретов в~исходном коде.
\end{itemize}

Результатом испытания считается наличие сформированных отчетов и~отсутствие критических блокеров, делающих невозможной приемку версии.
\end{hseExample}

\subsection{Испытание временных характеристик нагрузочным методом}

% Подсказка: подраздел нужен, если в ТЗ заданы временные характеристики;
% иначе удалите его.
\begin{hseExample}
Проверка подтверждает временные характеристики, заявленные в~техническом задании. Инструмент испытания~--- \hseFill{инструмент нагрузочного тестирования}.

Порядок проведения испытания:
\begin{enumerate}[label=\arabic*), leftmargin=1.25cm]
    \item подготовить рабочую станцию и~развернуть испытуемую программу в~тестовом контуре;
    \item выполнить нагрузочный прогон командой \hseFill{команда запуска прогона};
    \item выполнить ступенчатую серию прогонов с~параметрами \hseFill{число пользователей и длительность}; первые \hseFill{длительность периода прогрева} считать периодом прогрева и~исключать из~приемочных метрик;
    \item зафиксировать сводные показатели прогона: общее число запросов, число отказов, агрегированную интенсивность запросов, а~также перцентили времени отклика (p50, p95, p99);
    \item сохранить или приложить к~материалам испытаний отчет и~снимки экрана сводной статистики.
\end{enumerate}

Критерии успешного прохождения:
\begin{itemize}[nosep, leftmargin=1.25cm]
    \item агрегированная интенсивность запросов после периода прогрева составляет не~менее \hseFill{целевая интенсивность};
    \item доля успешных ответов после периода прогрева не~ниже \hseFill{целевая доля успешных ответов};
    \item отсутствуют необработанные исключения и~падения сервисов в~ходе прогона;
    \item сформирован отчет с~таблицей запросов, статистикой времени отклика и~графиками нагрузки.
\end{itemize}
\end{hseExample}

\subsection{Испытание трассировки требований и комплектности артефактов}

\begin{hseExample}
Для каждого функционального требования технического задания с~идентификатором вида \texttt{ТЗ-Ф-XX} должна быть указана автоматизированная проверка и~расположение соответствующего теста. Нефункциональные требования подтверждаются отдельными методами испытаний: \hseFill{методы проверки нефункциональных требований}. Дополнительно проверяется комплектность подтверждающих материалов испытаний, приложенных к~настоящему документу.
\end{hseExample}

\clearpage
% Страницы источников печатаются без нижнего штампа (как в реальных
% комплектах) — только номер листа и код документа сверху.
\pagestyle{appendix}

% ============================================================
% СПИСОК ИСПОЛЬЗОВАННЫХ ИСТОЧНИКОВ
% ============================================================
\section*{СПИСОК ИСПОЛЬЗОВАННЫХ ИСТОЧНИКОВ}
\addcontentsline{toc}{section}{СПИСОК ИСПОЛЬЗОВАННЫХ ИСТОЧНИКОВ}

% Подсказка: добавьте ПЕРЕД ГОСТами источники по технологиям вашего
% проекта в алфавитном порядке, по образцу:
% \bibitem{pytest}
% pytest Documentation [Электронный ресурс]. Режим доступа: \url{https://docs.pytest.org/}, свободный. (дата обращения: 14.03.2026).
\renewcommand{\refname}{}
\begin{thebibliography}{99}
\bibitem{gost19101}
ГОСТ 19.101-77: Виды программ и~программных документов. // Единая система программной документации.~--- М.: ИПК Издательство стандартов, 2001.

\bibitem{gost19103}
ГОСТ 19.103-77: Обозначения программ и~программных документов. // Единая система программной документации.~--- М.: ИПК Издательство стандартов, 2001.

\bibitem{gost19104}
ГОСТ 19.104-78: Основные надписи. // Единая система программной документации.~--- М.: ИПК Издательство стандартов, 2001.

\bibitem{gost19105}
ГОСТ 19.105-78: Общие требования к~программным документам. // Единая система программной документации.~--- М.: ИПК Издательство стандартов, 2001.

\bibitem{gost19106}
ГОСТ 19.106-78: Требования к~программным документам, выполненным печатным способом. // Единая система программной документации.~--- М.: ИПК Издательство стандартов, 2001.

\bibitem{gost19301}
ГОСТ 19.301-79: Программа и~методика испытаний. Требования к~содержанию и~оформлению. // Единая система программной документации.~--- М.: ИПК Издательство стандартов, 2001.

\bibitem{gost19604}
ГОСТ 19.604-78: Правила внесения изменений в~программные документы, выполненные печатным способом. // Единая система программной документации.~--- М.: ИПК Издательство стандартов, 2001.
\end{thebibliography}


\clearpage

% ============================================================
% ЛИСТ РЕГИСТРАЦИИ ИЗМЕНЕНИЙ
% ============================================================
\input{../common/change_log}

\end{document}
